\scanpage{13}
% archive foliation: 12
la désigne par $\unsure{\operatorname{Cor}|\operatorname{Cn}}_{X/Y}$, on vérifie aussitôt les conditions (i ter) : (iv ter).

\margin{Expliquer la relation de cet exemple avec le cas général, peut-on en \unsure{déduire} \uline{\unsure{galoisiennes}} \uline{\unsure{spéciales}} \uline{\unsure{par un autre}}}

\uline{Exemple II} Soit $C$ la catégorie des schémas \inserted{(dans un univers donné)}. Un morphisme fini et plat \inserted{$u:X\longrightarrow Y$} est dit de \uline{degré $r$} s'il correspond à un \deleted{\unsure{bon}} faisceau \deleted{cohérent de} \inserted{\unsure{qu. coh.}} \unsure{d'algèbres} sur $Y$, loc. libre de \deleted{degré} rang $r$ en $\forall$ point. Soit $\mathcal N$ l'ensemble des morphismes qui sont finis et plats et ont un \deleted{\unsure{un}} degré $r$. Les conditions (i) : (v) et (i bis) : (v bis) sont vérifiées, comme on voit \deleted{aussitôt} très facilement. \deleted{De} Plus, on peut définir de façon naturelle \uline{des} \deleted{applications} morphismes $\unsure{\operatorname{Cor}|\operatorname{Cn}}_{X/Y}$, satisfaisant les conditions (i) : (iv).

\uline{Exemple III \deleted{\unsure{Variantes}}} \inserted{(Variantes)} En géométrie analytique, on a une variante directe. En topologie, prenant les revêtements \unsure{non ramifiés} \deleted{de} \inserted{ayant un} degré, on trouve une variante analogue ; elle est très proche de l'exemple I, \deleted{\illegible{2}} \inserted{de façon} précise, soit $X$ un espace topologique \uline{fixé}, \deleted{prenons les} Alors
